
\documentclass[11pt,a4paper]{article}

%================================================
% Packages
%================================================

\usepackage{array}
\usepackage{float}
\usepackage{tabularx}
\usepackage{longtable}
\usepackage[utf8]{inputenc}
\usepackage[margin=1in]{geometry}
\usepackage{hyperref}
\usepackage{booktabs}
\usepackage{titlesec}
\usepackage{xcolor}
\usepackage{enumitem}
\usepackage{graphicx}


%================================================
% Color Scheme
%================================================

\definecolor{primary}{RGB}{30,41,59}
\definecolor{secondary}{RGB}{71,85,105}
\definecolor{accent}{RGB}{2,132,199}


\hypersetup{
    colorlinks=true,
    linkcolor=primary,
    urlcolor=accent
}


\titleformat{\section}
{\color{primary}\normalfont\Large\bfseries}
{\thesection}{1em}{}[\color{secondary}\titlerule]


\titleformat{\subsection}
{\color{secondary}\normalfont\large\bfseries}
{\thesubsection}{1em}{}


\setlist[itemize]{leftmargin=*}


%================================================
% Document
%================================================

\begin{document}



%================================================
% Title
%================================================


\begin{center}


{\LARGE\bfseries\color{primary}
Country Soccer Club System (CSCS)
}
\\[0.4cm]


{\large\color{secondary}
Database Design, Implementation, and Query Processing
}
\\[0.5cm]


{\large
COMP 353 -- Databases Project Report
}


\vspace{0.5cm}


{\large
Summer 2026
}


\end{center}



%------------------------------------------------
% Team Information
%------------------------------------------------

\section*{Administrative \& Team Information}


\begin{minipage}{0.45\textwidth}

\textbf{Course:} COMP 353 CC 1151\\
\textbf{Project:} CSCS Database Application\\
\textbf{Repository:}
\href{https://github.com/DatabasesSummerII2026/WarmUp_CSCS}
{WarmUpCSCS}

\end{minipage}
\hfill
\begin{minipage}{0.5\textwidth}

\begin{tabular}{ll}
\toprule
Name & Student ID\\
\midrule
Aaida Hossain & 40281796\\
Chunru Qiu & 40283505\\
Jingyuan Zhang & 40257596 \\
Yuqian Cai & 40187954\\

\bottomrule
\end{tabular}

\end{minipage}


%================================================
\section{Executive Summary}
%================================================


The Country Soccer Club System (CSCS) is a relational database
application developed to support the management activities of a
non-profit soccer organization.


The purpose of this project is to design and implement a database system
capable of managing club locations, personnel, members, family
relationships, teams, hobbies, financial payments, and FIFA game
participation records.


The system supports both major and minor club members using a unified
single-table inheritance model. Major members are members who are at
least 18 years old, while minor members (ages 4--17) require association with
registered family members or guardians. The database maintains membership
information, location movement history, family relationships, and player
development information over time.


The project follows a complete database development process, including
Entity-Relationship (E/R) modeling, relational schema conversion,
database implementation using MySQL, data population, and SQL query
development.


A total of fifteen relational tables were designed to represent the
exact specification requirements of the Country Soccer Club System. The design
applies database normalization principles, temporal composite primary key
structures, domain restrictions, and integrity rules to ensure data consistency.


The final database provides support for complex queries involving member
statistics, FIFA game participation, multi-installment payment analysis,
hobby tracking, and location-based administrative reports.



%================================================
\section{Introduction}
%================================================


The Country Soccer Club is a non-profit organization dedicated to
developing, promoting, and enhancing soccer activities for players aged
four years and above.


As the organization expands, managing information manually becomes
increasingly difficult. The club requires a centralized database system
to maintain accurate information about members, employees, locations,
teams, financial activities, and game participation.


The objective of this project is to develop a database application that
can support the operational requirements of the Country Soccer Club
System (CSCS).


The database must maintain information about:


\begin{itemize}


\item Club headquarters (Head) and branch locations, including manager assignments.


\item Personnel working at different locations over time with specific roles and mandates.


\item Major and minor club members, including registration age enforcement.


\item Family members and guardian relationships.


\item Member hobbies and personal information.


\item Membership payments, installment tracking, and donation management.


\item Soccer teams and FIFA game participation history.


\end{itemize}



The system was designed according to the requirements provided in the
COMP 353 Database course project specification.


The development process consists of several major stages:


\begin{itemize}


\item Requirement analysis and identification of entities.


\item Entity-Relationship model design.


\item Conversion from E/R model to relational schema.


\item MySQL database implementation.


\item Data insertion and validation.


\item SQL query formulation and testing.


\end{itemize}



The final system provides a structured and reliable database solution
that allows the Country Soccer Club to efficiently manage its daily
operations and support future expansion.


%================================================
\section{System Requirement Analysis}
%================================================


The Country Soccer Club System requirements were analyzed to identify
the major entities, relationships, constraints, and operational
requirements of the organization.



%------------------------------------------------
\subsection{Location Management}


The system maintains information about all club locations.


A location can represent:

\begin{itemize}

\item The main headquarters (\texttt{Head}).

\item A branch location (\texttt{Branch}).

\end{itemize}


Each location contains a unique ID, name, address details, contact phone,
website URL, maximum member capacity, and a reference to the managing personnel.
A personnel member can manage at most one location at a time.



%------------------------------------------------
\subsection{Personnel Management}


The system maintains information about every employee or volunteer
working for the club.


Personnel information includes personal demographic data, contact info,
unique SSN, Medicare number, role, and mandate (\texttt{Volunteer} or \texttt{Salaried}).


Personnel roles include:
\texttt{Administrator}, \texttt{Captain}, \texttt{Coach}, \texttt{Assistant Coach}, or \texttt{Other}.


A personnel member can work at different locations during different time
periods. Therefore, a temporal relation (\texttt{Personnel\_Assignment}) with
\texttt{startDate} and \texttt{endDate} is used to track assignment history.



%------------------------------------------------
\subsection{Club Member Management}


The system maintains information about all registered club members using a globally
unique, auto-incremented membership number.


Members are classified into:


\begin{itemize}

\item \textbf{Major members}: Age 18 or above.

\item \textbf{Minor members}: Age between 4 and 17.

\end{itemize}


The database enforces that a member must be at least 4 years old at the time of
registration (\texttt{registrationDate} $\ge$ \texttt{DOB + 4 years}). Physical metrics
(height and weight), SSN, Medicare number, and location movement history are tracked.



%------------------------------------------------
\subsection{Family Member Management}


Minor members must be associated with at least one registered family
member or guardian.


The system maintains family demographics, contact details, and a temporal
relationship (\texttt{FamilyRelationship}) capturing guardian dynamic associations
(\texttt{Father}, \texttt{Mother}, \texttt{Grandfather}, \texttt{Grandmother}, \texttt{Tutor}, \texttt{Partner}, \texttt{Friend}, \texttt{Other}).



%------------------------------------------------
\subsection{Financial Management}


The system tracks all membership payments per member and membership year.


Payments can be paid in up to four installments per year. The required annual fee structure is:


\begin{table}[H]

\centering


\begin{tabular}{cc}


\toprule

Member Type & Standard Annual Fee \\


\midrule


Minor Member & \$100 \\


Major Member & \$200 \\


\bottomrule


\end{tabular}


\caption{Annual Membership Fees}


\end{table}



Any total amount paid by a member within a membership year that exceeds the required fee is recorded as a donation. If a member's fee for the preceding year was not fully paid, the member's status is considered inactive.



%------------------------------------------------
\subsection{Team and FIFA Game Management}


Teams are created at specific locations and categorized by gender (\texttt{Boys} or \texttt{Girls}).
Members can join teams over specified date intervals. The system tracks official FIFA games,
including game date, location, opponent team, final score, and player participation.





%================================================
\section{Database Design Decisions}
%================================================


\subsection{Temporal Relationships \& Composite Primary Keys}

To satisfy requirements where entities change locations or roles over time, temporal attributes (\texttt{startDate} and \texttt{endDate}) are embedded into historical tables. The composite primary key includes \texttt{startDate} (e.g., \texttt{(memberID, locationID, startDate)}) to allow a member or personnel to be re-assigned to the same location at different times. A \texttt{NULL} value in \texttt{endDate} signifies an active association.


\subsection{Single-Table Inheritance for Members}

Club members (\texttt{Major} and \texttt{Minor}) are implemented in a single table (\texttt{ClubMembers}) using a \texttt{memberType} discriminator column. This simplifies foreign key references in payment and game participation tables while preserving type-specific constraints via check constraints.


\subsection{Normalization}

All relations satisfy 3NF (Third Normal Form):
\begin{itemize}
\item Every non-key attribute is fully functionally dependent on the primary key.
\item Transitive dependencies are eliminated (e.g., location manager info is referenced via \texttt{managerID} rather than embedded in \texttt{Locations}).
\end{itemize}





%================================================
\section{Relational Schema Design}
%================================================


The final CSCS relational schema consists of fifteen normalized relations as defined below. Primary keys are underlined.


%------------------------------------------------
\subsection{Locations}

\begin{table}[H]
\centering
\begin{tabular}{lll}
\toprule
Attribute & Data Type & Constraint \\
\midrule
\underline{locationID} & INT & Primary Key, Auto-Increment \\
name & VARCHAR(100) & NOT NULL \\
type & VARCHAR(20) & CHECK (\texttt{Head}, \texttt{Branch}) \\
address & VARCHAR(255) & \\
city & VARCHAR(100) & \\
province & VARCHAR(100) & \\
postalCode & VARCHAR(20) & \\
phone & VARCHAR(20) & \\
webAddress & VARCHAR(255) & \\
capacity & INT & \\
managerID & INT & Foreign Key $\rightarrow$ Personnel(personnelID), UNIQUE \\
\bottomrule
\end{tabular}
\caption{Locations Relation}
\end{table}


%------------------------------------------------
\subsection{Personnel}

\begin{table}[H]
\centering
\begin{tabular}{lll}
\toprule
Attribute & Data Type & Constraint \\
\midrule
\underline{personnelID} & INT & Primary Key, Auto-Increment \\
firstName & VARCHAR(100) & NOT NULL \\
lastName & VARCHAR(100) & NOT NULL \\
DOB & DATE & NOT NULL \\
SSN & VARCHAR(20) & UNIQUE, NOT NULL \\
MedicareNo & VARCHAR(50) & UNIQUE \\
phone & VARCHAR(20) & \\
email & VARCHAR(100) & \\
address & VARCHAR(255) & \\
city & VARCHAR(100) & \\
province & VARCHAR(100) & \\
postalCode & VARCHAR(20) & \\
role & VARCHAR(50) & CHECK (\texttt{Administrator}, \texttt{Captain}, \texttt{Coach}, \dots) \\
mandate & VARCHAR(50) & CHECK (\texttt{Volunteer}, \texttt{Salaried}) \\
\bottomrule
\end{tabular}
\caption{Personnel Relation}
\end{table}


%------------------------------------------------
\subsection{ClubMembers}

\begin{table}[H]
\centering
\begin{tabular}{lll}
\toprule
Attribute & Data Type & Constraint \\
\midrule
\underline{memberID} & INT & Primary Key, Auto-Increment \\
firstName & VARCHAR(100) & NOT NULL \\
lastName & VARCHAR(100) & NOT NULL \\
DOB & DATE & NOT NULL \\
height & DECIMAL(5,2) & \\
weight & DECIMAL(5,2) & \\
SSN & VARCHAR(20) & UNIQUE, NOT NULL \\
MedicareNo & VARCHAR(50) & UNIQUE \\
phone & VARCHAR(20) & \\
address & VARCHAR(255) & \\
city & VARCHAR(100) & \\
province & VARCHAR(100) & \\
postalCode & VARCHAR(20) & \\
registrationDate & DATE & NOT NULL, CHECK ($\ge$ DOB + 4 years) \\
memberType & VARCHAR(10) & CHECK (\texttt{Major}, \texttt{Minor}) \\
\bottomrule
\end{tabular}
\caption{ClubMembers Relation}
\end{table}


%------------------------------------------------
\subsection{FamilyMembers}

\begin{table}[H]
\centering
\begin{tabular}{lll}
\toprule
Attribute & Data Type & Constraint \\
\midrule
\underline{familyID} & INT & Primary Key, Auto-Increment \\
firstName & VARCHAR(100) & NOT NULL \\
lastName & VARCHAR(100) & NOT NULL \\
DOB & DATE & \\
SSN & VARCHAR(20) & UNIQUE, NOT NULL \\
MedicareNo & VARCHAR(50) & UNIQUE \\
phone & VARCHAR(20) & \\
email & VARCHAR(100) & \\
address & VARCHAR(255) & \\
city & VARCHAR(100) & \\
province & VARCHAR(100) & \\
postalCode & VARCHAR(20) & \\
\bottomrule
\end{tabular}
\caption{FamilyMembers Relation}
\end{table}


%------------------------------------------------
\subsection{Hobbies}

\begin{table}[H]
\centering
\begin{tabular}{lll}
\toprule
Attribute & Data Type & Constraint \\
\midrule
\underline{hobbyID} & INT & Primary Key, Auto-Increment \\
hobbyName & VARCHAR(100) & UNIQUE, NOT NULL \\
\bottomrule
\end{tabular}
\caption{Hobbies Relation}
\end{table}


%------------------------------------------------
\subsection{Personnel\_Assignment}

\begin{table}[H]
\centering
\begin{tabular}{lll}
\toprule
Attribute & Data Type & Constraint \\
\midrule
\underline{personnelID} & INT & Foreign Key $\rightarrow$ Personnel(personnelID) \\
\underline{locationID} & INT & Foreign Key $\rightarrow$ Locations(locationID) \\
\underline{startDate} & DATE & Primary Key Part \\
endDate & DATE & Nullable (NULL = Active) \\
\bottomrule
\end{tabular}
\caption{Personnel Assignment Relation}
\end{table}


%------------------------------------------------
\subsection{Member\_Location\_History}

\begin{table}[H]
\centering
\begin{tabular}{lll}
\toprule
Attribute & Data Type & Constraint \\
\midrule
\underline{memberID} & INT & Foreign Key $\rightarrow$ ClubMembers(memberID) \\
\underline{locationID} & INT & Foreign Key $\rightarrow$ Locations(locationID) \\
\underline{startDate} & DATE & Primary Key Part \\
endDate & DATE & Nullable (NULL = Active) \\
\bottomrule
\end{tabular}
\caption{Member Location History Relation}
\end{table}


%------------------------------------------------
\subsection{Family\_Location\_History}

\begin{table}[H]
\centering
\begin{tabular}{lll}
\toprule
Attribute & Data Type & Constraint \\
\midrule
\underline{familyID} & INT & Foreign Key $\rightarrow$ FamilyMembers(familyID) \\
\underline{locationID} & INT & Foreign Key $\rightarrow$ Locations(locationID) \\
\underline{startDate} & DATE & Primary Key Part \\
endDate & DATE & Nullable (NULL = Active) \\
\bottomrule
\end{tabular}
\caption{Family Location History Relation}
\end{table}


%------------------------------------------------
\subsection{FamilyRelationship}

\begin{table}[H]
\centering
\begin{tabular}{lll}
\toprule
Attribute & Data Type & Constraint \\
\midrule
\underline{familyID} & INT & Foreign Key $\rightarrow$ FamilyMembers(familyID) \\
\underline{memberID} & INT & Foreign Key $\rightarrow$ ClubMembers(memberID) \\
\underline{startDate} & DATE & Primary Key Part \\
endDate & DATE & Nullable \\
relationshipType & VARCHAR(50) & CHECK (\texttt{Father}, \texttt{Mother}, \texttt{Tutor}, \dots) \\
\bottomrule
\end{tabular}
\caption{Family Relationship Relation}
\end{table}


%------------------------------------------------
\subsection{Member\_Hobby}

\begin{table}[H]
\centering
\begin{tabular}{lll}
\toprule
Attribute & Data Type & Constraint \\
\midrule
\underline{memberID} & INT & Foreign Key $\rightarrow$ ClubMembers(memberID) \\
\underline{hobbyID} & INT & Foreign Key $\rightarrow$ Hobbies(hobbyID) \\
\bottomrule
\end{tabular}
\caption{Member Hobby Relation}
\end{table}


%------------------------------------------------
\subsection{Payments}

\begin{table}[H]
\centering
\begin{tabular}{lll}
\toprule
Attribute & Data Type & Constraint \\
\midrule
\underline{paymentID} & INT & Primary Key, Auto-Increment \\
memberID & INT & Foreign Key $\rightarrow$ ClubMembers(memberID) \\
paymentDate & DATE & NOT NULL \\
amount & DECIMAL(10,2) & NOT NULL \\
method & VARCHAR(20) & CHECK (\texttt{Cash}, \texttt{Debit}, \texttt{Credit Card}) \\
membershipYear & INT & NOT NULL \\
installmentNumber & INT & CHECK (BETWEEN 1 AND 4) \\
\bottomrule
\end{tabular}
\caption{Payments Relation (UNIQUE on \texttt{(memberID, membershipYear, installmentNumber)})}
\end{table}


%------------------------------------------------
\subsection{Teams}

\begin{table}[H]
\centering
\begin{tabular}{lll}
\toprule
Attribute & Data Type & Constraint \\
\midrule
\underline{teamID} & INT & Primary Key, Auto-Increment \\
teamName & VARCHAR(100) & NOT NULL \\
gender & VARCHAR(10) & CHECK (\texttt{Boys}, \texttt{Girls}) \\
locationID & INT & Foreign Key $\rightarrow$ Locations(locationID) \\
\bottomrule
\end{tabular}
\caption{Teams Relation}
\end{table}


%------------------------------------------------
\subsection{Team\_Members}

\begin{table}[H]
\centering
\begin{tabular}{lll}
\toprule
Attribute & Data Type & Constraint \\
\midrule
\underline{teamID} & INT & Foreign Key $\rightarrow$ Teams(teamID) \\
\underline{memberID} & INT & Foreign Key $\rightarrow$ ClubMembers(memberID) \\
\underline{startDate} & DATE & Primary Key Part \\
endDate & DATE & Nullable \\
\bottomrule
\end{tabular}
\caption{Team Members Relation}
\end{table}


%------------------------------------------------
\subsection{FIFA\_Games}

\begin{table}[H]
\centering
\begin{tabular}{lll}
\toprule
Attribute & Data Type & Constraint \\
\midrule
\underline{gameID} & INT & Primary Key, Auto-Increment \\
gameDate & DATE & NOT NULL \\
locationID & INT & Foreign Key $\rightarrow$ Locations(locationID) \\
opponentTeam & VARCHAR(100) & NOT NULL \\
score & VARCHAR(20) & \\
\bottomrule
\end{tabular}
\caption{FIFA Games Relation}
\end{table}


%------------------------------------------------
\subsection{Game\_Participation}

\begin{table}[H]
\centering
\begin{tabular}{lll}
\toprule
Attribute & Data Type & Constraint \\
\midrule
\underline{memberID} & INT & Foreign Key $\rightarrow$ ClubMembers(memberID) \\
\underline{gameID} & INT & Foreign Key $\rightarrow$ FIFA\_Games(gameID) \\
teamID & INT & Foreign Key $\rightarrow$ Teams(teamID) \\
\bottomrule
\end{tabular}
\caption{Game Participation Relation}
\end{table}





%================================================
\section{SQL Query Implementation}
%================================================


The CSCS database supports required operational and reporting queries. Below are the key query implementations updated for the target schema.


%------------------------------------------------
\subsection{Query I: Location Information Report}

Retrieves information for every location, including manager details, count of personnel assigned, active members, and game participants.

\begin{verbatim}
SELECT 
    L.locationID,
    L.name AS LocationName,
    L.type,
    CONCAT(P.firstName, ' ', P.lastName) AS ManagerName,
    COUNT(DISTINCT PA.personnelID) AS TotalPersonnel,
    COUNT(DISTINCT MLH.memberID) AS TotalMembers
FROM Locations L
LEFT JOIN Personnel P ON L.managerID = P.personnelID
LEFT JOIN Personnel_Assignment PA ON L.locationID = PA.locationID 
    AND PA.endDate IS NULL
LEFT JOIN Member_Location_History MLH ON L.locationID = MLH.locationID 
    AND MLH.endDate IS NULL
GROUP BY L.locationID, L.name, L.type, ManagerName
ORDER BY TotalMembers DESC;
\end{verbatim}


%------------------------------------------------
\subsection{Query II: Major Members Who Played FIFA Games}

\begin{verbatim}
SELECT 
    CM.memberID,
    CM.firstName,
    CM.lastName,
    TIMESTAMPDIFF(YEAR, CM.DOB, CURDATE()) AS Age,
    COUNT(GP.gameID) AS GamesPlayed
FROM ClubMembers CM
JOIN Game_Participation GP ON CM.memberID = GP.memberID
WHERE CM.memberType = 'Major'
GROUP BY CM.memberID, CM.firstName, CM.lastName, CM.DOB
ORDER BY GamesPlayed DESC;
\end{verbatim}


%------------------------------------------------
\subsection{Query III: Members With Four Or More Hobbies}

\begin{verbatim}
SELECT 
    CM.memberID,
    CM.firstName,
    CM.lastName,
    COUNT(MH.hobbyID) AS HobbyCount
FROM ClubMembers CM
JOIN Member_Hobby MH ON CM.memberID = MH.memberID
GROUP BY CM.memberID, CM.firstName, CM.lastName
HAVING COUNT(MH.hobbyID) >= 4
ORDER BY HobbyCount DESC;
\end{verbatim}


%------------------------------------------------
\subsection{Query IV: Major Members Without FIFA Participation}

\begin{verbatim}
SELECT 
    CM.memberID,
    CM.firstName,
    CM.lastName
FROM ClubMembers CM
LEFT JOIN Game_Participation GP ON CM.memberID = GP.memberID
WHERE CM.memberType = 'Major' AND GP.gameID IS NULL;
\end{verbatim}


%------------------------------------------------
\subsection{Query V: Member Age Statistics}

\begin{verbatim}
SELECT 
    TIMESTAMPDIFF(YEAR, DOB, CURDATE()) AS Age,
    COUNT(memberID) AS TotalMembers
FROM ClubMembers
GROUP BY Age
ORDER BY Age DESC;
\end{verbatim}


%------------------------------------------------
\subsection{Query VI: Major Members Who Are Also Family Members}

\begin{verbatim}
SELECT 
    CM.memberID AS ParentMemberID,
    CM.firstName AS ParentFirstName,
    CM.lastName AS ParentLastName,
    Child.memberID AS ChildMemberID,
    Child.firstName AS ChildFirstName,
    Child.lastName AS ChildLastName,
    FR.relationshipType
FROM ClubMembers CM
JOIN FamilyMembers FM ON CM.SSN = FM.SSN
JOIN FamilyRelationship FR ON FM.familyID = FR.familyID
JOIN ClubMembers Child ON FR.memberID = Child.memberID
WHERE CM.memberType = 'Major' AND FR.endDate IS NULL;
\end{verbatim}


%------------------------------------------------
\subsection{Query VII: Membership Fees and Donations Report}

Calculates fee contributions (\$200 standard fee for Major members) and donations aggregated from total payments across installments between 2023 and 2025.

\begin{verbatim}
SELECT 
    MemberTotals.membershipYear,
    SUM(LEAST(MemberTotals.TotalPaid, 200)) AS TotalFees,
    SUM(GREATEST(MemberTotals.TotalPaid - 200, 0)) AS TotalDonations
FROM (
    SELECT 
        P.memberID,
        P.membershipYear,
        SUM(P.amount) AS TotalPaid
    FROM Payments P
    JOIN ClubMembers CM ON P.memberID = CM.memberID
    WHERE CM.memberType = 'Major' 
      AND P.membershipYear BETWEEN 2023 AND 2025
    GROUP BY P.memberID, P.membershipYear
) AS MemberTotals
GROUP BY MemberTotals.membershipYear;
\end{verbatim}


%------------------------------------------------
\subsection{Query VIII: Members Participating in Four or More FIFA Games}

\begin{verbatim}
SELECT 
    CM.memberID,
    CM.firstName,
    CM.lastName,
    CM.memberType,
    COUNT(GP.gameID) AS TotalGames
FROM ClubMembers CM
JOIN Game_Participation GP ON CM.memberID = GP.memberID
GROUP BY CM.memberID, CM.firstName, CM.lastName, CM.memberType
HAVING COUNT(GP.gameID) >= 4
ORDER BY TotalGames DESC;
\end{verbatim}





%================================================
\section{Database Relation Tuple Count}
%================================================


\begin{table}[H]
\centering
\begin{tabular}{lc}
\toprule
Relation Name & Tuple Count \\
\midrule
Locations & 10 \\
Personnel & 10 \\
Personnel\_Assignment & 10 \\
ClubMembers & 10 \\
Member\_Location\_History & 10 \\
FamilyMembers & 10 \\
Family\_Location\_History & 10 \\
FamilyRelationship & 10 \\
Hobbies & 10 \\
Member\_Hobby & 10 \\
Payments & 10 \\
Teams & 10 \\
Team\_Members & 10 \\
FIFA\_Games & 10 \\
Game\_Participation & 10 \\
\bottomrule
\end{tabular}
\caption{Database Population Summary}
\end{table}





%================================================
\section{Conclusion}
%================================================


The Country Soccer Club System (CSCS) database was fully designed and implemented on MySQL 8.0. The schema accurately models physical location operations, personnel mandates, player demographics, multi-installment payments, and competitive play history while strictly maintaining relational integrity and normalization standards.





\begin{thebibliography}{9}

\bibitem{mysql}
MySQL 8.0 Reference Manual,
\url{https://dev.mysql.com/doc/refman/8.0/en/}

\bibitem{elmasri}
R. Elmasri and S. Navathe,
\textit{Fundamentals of Database Systems}, 7th ed., Pearson.

\bibitem{silberschatz}
A. Silberschatz, H. Korth, and S. Sudarshan,
\textit{Database System Concepts}, 7th ed., McGraw-Hill.

\end{thebibliography}



\end{document}
